\chapter{Related Work}


\section{Literature Overview and Key Issues}

Automatic grading systems have long been used in programming education. Their typical logic is "compilation + test cases + score by pass rate." The advantages are clear: high speed, unified rules, and suitability for batch submissions. In larger classes, this approach indeed reduces repetitive instructor workload.

However, traditional automatic grading also has notable limitations in practical teaching. First, it focuses mainly on functional correctness and provides limited coverage of readability, structural design, and coding standards. Second, feedback is usually brief, often limited to pass/fail outcomes and difficult to interpret. Third, when code has minor syntax errors but demonstrates correct core reasoning, systems may still assign very low scores or zero, which is not always consistent with how instructors evaluate student thinking in manual review. In other words, test-driven scoring is efficient but still leaves room for improvement in pedagogical feedback quality.

\paragraph{Application of Large Language Models in Code Feedback}
In recent years, large language models have shown strong performance in code generation, explanation, and error analysis, creating a new path for automated explanatory feedback. Compared with traditional automatic grading, LLMs can produce richer natural-language explanations, such as identifying potential logic issues, explaining compilation errors, and suggesting revisions.

From a modeling and evaluation perspective, the Codex study emphasizes functional correctness as the core target and evaluates generated code with pass@k, that is, whether at least one correct solution appears within k sampled attempts under unit tests, rather than relying on surface-form similarity alone \cite{ref5_chen2021}. The same work also reports weak correlation between BLEU and true functional correctness. Austin et al. report a similar pattern in program synthesis: larger models generally improve solve rates, while sampling strategy strongly affects outcomes; higher-temperature sampling can be more effective when multiple attempts are allowed \cite{ref6_austin2021}. These observations directly support the design choice in this thesis to prioritize executable evidence, including runtime behavior and hidden-test outcomes, over text-level similarity.

Looking at concrete experimental settings, Chen et al. report pass@k under different sampling budgets on HumanEval and emphasize that temperature should be tuned with respect to k rather than fixed globally \cite{ref5_chen2021}. Austin et al. also compare greedy decoding with temperature-based sampling in program synthesis and show that, when multiple attempts are allowed, diverse sampling often finds executable correct programs more reliably than a single highest-probability output \cite{ref6_austin2021}. This is directly relevant to teaching systems: feedback should preserve executable evidence and reproducible outcomes, not depend only on the fluency of one generated response.

However, LLMs are not directly ready for classroom use. During investigation and development, two common issues were observed: suggestions may go beyond current course progress, and responses may be overly generic or inconsistent with course terminology. Therefore, in teaching scenarios, LLMs are better used as a feedback enhancement module rather than a full replacement for instructional grading rules.

\paragraph{Role of RAG in Teaching Feedback}
To address the problem of feedback drifting away from course context, Retrieval-Augmented Generation (RAG) provides a feasible approach. The basic process is to retrieve relevant content from course documents first, and then generate feedback using both retrieved materials and submission context. Compared with direct generation, RAG has two key advantages: stronger alignment with course content (fewer out-of-scope suggestions) and improved interpretability through material grounding.

In the original RAG framework, the generator combines parametric memory with non-parametric external memory by conditioning on retrieved documents and marginalizing over top-k candidates during generation \cite{ref3_lewis2020}. The paper further distinguishes RAG-Sequence and RAG-Token, where the token-level variant can switch evidence across tokens. For teaching feedback, this mechanism is important because it keeps responses grounded in editable course materials while still benefiting from the expressive power of large models, reducing hallucination and terminology drift.

In this project, lecture notes and review materials are included in the retrieval corpus. Feedback generation references retrieved passages to keep responses close to classroom context rather than producing generic suggestions.

\paragraph{Academic Integrity and Fairness}
In programming courses, an automated system should be not only fast but also fair. Fairness is mainly reflected in two points: applying consistent rules to all students and ensuring scores reflect actual learning effort. If the system cannot detect direct template submissions or disguised blank submissions, unreasonable high scores may occur and undermine fairness. Conversely, if syntax errors always result in immediate failure, submissions with meaningful reasoning may be overly penalized.

To address this, the project adopts a separated final/static scoring strategy and introduces blank-submission detection. This helps balance result correctness and process evaluation, reducing unfairness caused by scoring distortion.

\section{Chapter Summary}

Overall, existing studies and engineering practice have built a foundation in efficiency, interpretability, and course alignment, but a fully deployable solution for undergraduate teaching is still limited. This thesis does not focus on introducing a new algorithm; instead, it integrates automatic grading, rule-based scoring, RAG feedback, and instructor management into a deployable, operable, and iterative platform for real course workflows.

